% TOP_PROBLEM: {"id": 3041, "title": "Even vs. odd latin squares", "category_id": 2, "category": {"id": 2, "name": "combinatorics", "display_name": "Combinatorics", "description": "Counting problems, graph theory, discrete structures.", "slug": "combinatorics", "order_index": 2, "created_at": "2026-07-31T15:26:25.671Z"}, "status": "open", "problem_number": "OPG-636", "set_id": 12}
% FIRST_POSTED: 2026-10-01
% ATTEMPT_STATUS: unresolved
% UnsolvedMath ID 3041; source catalogue code OPG-636
% !TeX program = lualatex
% Research attempt by GPT-6 Astra Ultra; no claim of a new complete solution.
\documentclass[11pt,a4paper]{article}
\usepackage[margin=25mm]{geometry}
\usepackage{fontspec}
\setmainfont{lmroman10-regular.otf}[BoldFont=lmroman10-bold.otf,
 ItalicFont=lmroman10-italic.otf,BoldItalicFont=lmroman10-bolditalic.otf]
\usepackage{amsmath,amssymb,mathtools}
\usepackage{unicode-math}
\setmathfont{latinmodern-math.otf}
\AtBeginDocument{\let\setminus\smallsetminus}
\usepackage{xurl,hyperref}
\hypersetup{colorlinks=true,urlcolor=blue}
\setcounter{secnumdepth}{0}
\setlength{\parindent}{0pt}
\setlength{\parskip}{5pt}
\setlength{\emergencystretch}{3em}
\begin{document}
\section*{Even vs. odd latin squares}\label{3041}
{\small UnsolvedMath ID 3041 · OPG-636 · Combinatorics · OpenGarden}
\subsection{Definitions and mathematical statement}
For a positive integer $n$, a Latin square of order $n$ is an array
$L=(L_{ij})_{1\le i,j\le n}$ with entries in $[n]=\{1,\ldots,n\}$, in which
each symbol occurs exactly once in each row and exactly once in each column.
The rows, columns and symbols are labelled: different labelled arrays are
counted separately, even if one can be obtained from another by relabelling.

Each row determines a permutation $\rho_i:j\mapsto L_{ij}$ of $[n]$, and
each column determines a permutation $\kappa_j:i\mapsto L_{ij}$.
For a permutation $\sigma$, let $\operatorname{sgn}(\sigma)$ be $+1$ when
its inversion count is even and $-1$ when that count is odd. Define
\[
 \varepsilon(L)=\prod_{i=1}^n\operatorname{sgn}(\rho_i)
                 \prod_{j=1}^n\operatorname{sgn}(\kappa_j).
\]
A Latin square is even if $\varepsilon(L)=1$, and odd otherwise. Write
$E_n$ and $O_n$ for the respective numbers of labelled squares.

The conjecture asserts $E_n\ne O_n$ for every positive even integer $n$,
or equivalently
\[
 \sum_{L\text{ Latin of order }n}\varepsilon(L)\ne0.
\]
The assertion is nonvanishing, rather than a prescribed sign or a lower
bound on the difference. It concerns the product of both row and column
signs, not just one family. No normalization of the first row or first
column is imposed in these counts; any reduction to normalized squares
must account for the effect of relabelling on parity.
\subsection{Short English statement}
For every even order, do the numbers of even and odd labelled Latin squares differ?
\subsection{Research attempt}
\paragraph{Scope and process.}
This is a GPT-6 Astra Ultra research attempt dated 1 October 2026, using
primary-source browsing and elementary mathematical checks. The canonical
definition above is retained. Results below are restricted results or
reductions, not a claim of a new solution to the entire catalogue question.
No formal proof assistant or external mathematical referee checked this text.


\paragraph{Literature and aim.}
This is the Alon--Tarsi conjecture with both row and column signs.
Drisko's primary paper [R1] proves nonvanishing for order $p+1$ when
$p$ is an odd prime. We do not extrapolate this to arbitrary even
orders. The present attempt derives a coefficient reformulation,
checks the parity normalization, and verifies orders two and four.

\paragraph{Relabelling and the odd-order obstruction.}
Permuting the rows by $\alpha$ merely reorders their signs, while
composing each column permutation with the row permutation or its
inverse multiplies its sign by $\operatorname{sgn}(\alpha)$.
Thus the square's sign changes by $\operatorname{sgn}(\alpha)^n$.
The same statement holds for a permutation of columns. A symbol
permutation $\beta$ composes every row and column permutation with
$\beta$, producing the factor $\operatorname{sgn}(\beta)^{2n}=1$.
For odd $n>1$, a fixed transposition of two rows is therefore a
fixed-point-free sign-reversing involution: distinct Latin rows
cannot be equal. It follows that $E_n=O_n$ in odd orders. For even
orders, all these relabellings preserve sign, explaining why this
simple cancellation mechanism disappears but not proving nonvanishing.

\paragraph{An exact polynomial coefficient.}
Use independent variables $x_{ij}$ and define
\[
 P_n(X)=
 \prod_{i=1}^n\prod_{j<k}(x_{ik}-x_{ij})\,
 \prod_{j=1}^n\prod_{i<k}(x_{kj}-x_{ij}).
\]
The Vandermonde expansion in each row assigns exponents
$0,\ldots,n-1$ bijectively to its cells, with coefficient the sign
of that exponent permutation. The column factors do the same.
To contribute to $M=\prod_{i,j}x_{ij}^{n-1}$, the column exponent
at each cell must be $n-1$ minus its row exponent. Thus the row
exponents plus one form a Latin square: they are a permutation in
each row by construction, and in each column because the complementary
column exponents are a permutation. Conversely every Latin square
gives exactly one such choice of row and column expansion terms.

Reversal of $0,\ldots,n-1$ has sign $(-1)^{\binom n2}$.
Consequently each column exponent permutation contributes that factor
times the sign of its Latin column. Multiplying over columns gives
the exact identity
\[
 [M]P_n(X)=(-1)^{n\binom n2}(E_n-O_n).
\]
For even $n$, the factor is $+1$. The target is therefore precisely
nonvanishing of this specified coefficient in every even order. A
polynomial being nonzero is not sufficient: individual coefficients
can cancel in its expansion, and that is the issue isolated here.

\paragraph{Small exact checks and their reproducibility.}
For $n=2$ the two labelled squares both have sign $+1$, so the
difference is two. For $n=4$, an exhaustive Python enumeration in
this attempt gave $E_4=576$, $O_4=0$. The enumeration builds an array
row by row, trying each of the $24$ permutations as a row and rejecting
a row whenever some column repeats a symbol already used. At depth
four it multiplies the signs of the four rows and four columns,
where each sign is $(-1)$ to its inversion count. Every labelled
Latin square occurs once because its ordered rows are unique, and
every accepted array is Latin because each column has four distinct
entries from a four-symbol set. This explains completeness of the
finite check; it is not a check of an asymptotic or all-order claim.
The exact enumeration is reproducible with
\texttt{scripts/check\_attempt\_batch\_algebra\_2026\_10\_01.py}.

\paragraph{Normalization audit and remaining gap.}
For even $n$, reduction of the first row and first column preserves
sign. Each reduced square accounts for $n!(n-1)!$ labelled squares:
choose the ordering of columns to set the first row, then the ordering
of the other rows to set the first column. Equivalently, normalize
these in reverse; there is exactly one choice at each step. Therefore
the signed labelled count is this positive factor times the signed
reduced count. Normalization cannot create the desired imbalance.
The first remaining gap is an all-order argument that the displayed
coefficient does not cancel. None of the sign transformations or
finite enumerations provides such an argument.

\paragraph{Final status.}
\textbf{UNRESOLVED.} The coefficient identity, symmetry reductions,
and small cases are checked; the general even-order conjecture is
not proved here.

\subsection{Sources}
\begin{itemize}
\item[S1] ULAM.AI and the UnsolvedMath Contributors, supplied problems.json, record 3041 (OPG-636), OpenGarden collection. Catalogue identity and source transcription; CC BY 4.0.
\url{https://huggingface.co/datasets/ulamai/UnsolvedMath}.
\item[S2] Open Problem Garden, Even vs. odd latin squares. Original problem and discussion; the mathematical formulation below is expanded from the supplied source record.
\url{http://www.openproblemgarden.org/op/even_vs_odd_latin_squares}.

\item[R1] Arthur A. Drisko, \emph{On the Number of Even and Odd Latin
Squares of Order $p+1$}, Advances in Mathematics 128 (1997), 20--35.
Primary publisher abstract checked for the $p+1$ theorem.
\url{https://www.sciencedirect.com/science/article/pii/S0001870897916236}.

\end{itemize}
\end{document}
